\subsection{Parameter hypotheses and their uses}\label{sec:hypotheses}

The labels used in the parameter verification below have the following mathematical meaning. For either keyhole,
its positive parameters satisfy $0<W_1<W_2\le1/6$, $0<\rho_1<1$ and the exact relations
\eqref{eq:K123}; its budgets $\delta_{\mathrm{rel}},\delta_{\mathrm{small}}$ are those of
Theorem~\ref{thm:kh-bounds}. Subscript $B$ denotes the corner keyhole.

\begin{description}[leftmargin=4em,labelwidth=3.5em]
\item[(H0)] $k\in\Z$, $k\ge10$, $0<\eps_c\le1/100$ (the parameter regime).
\item[(H1)] $0<\lambda_2<\lambda_a<1$ and $0<\mu<1$ (the comparator schedules).
\item[(H2)] $0<\rho_b<a$, equivalently $r_+>1$, where $\rho_b=R_s^{-2}$ and $a=r_+^2\rho_b$
(the enlarged source domain).
\item[(H3)--(H5)] $K_C$ has input radius $r_a=\sqrt{a/\pi}$ and satisfies the positive-parameter conditions and
\eqref{eq:K123} (Theorem~\ref{thm:kh-bounds} for the comparator).
\item[(H6)--(H7)] The two schedule inequalities displayed in Theorem~\ref{thm:comparator} hold for every
$u\in[0,1]$ (stack containment).
\item[(H8)] $\pi d^2\le T(1-\mu-\lambda_2/2)$ and
$\pi d^2+\lambda_2T/2+\mu T\le O(-w_G)$ (the remaining comparator branches).
\item[(H9)] $c_s\pi r_h^2+\frac32\lambda_2g\le g$. This is an additional implementation guard, also verified
below; the proof of Theorem~\ref{thm:comparator} uses (H6)--(H8) directly.
\item[(H9b)] $3-9a>\delta$, where $\delta=w_2/4$ (separation of the group windows from the marked-fibre spirals).
\item[(H10)] The constants and integer shifts are those of Construction~\ref{cons:twists}; its four $v$-windows
are pairwise disjoint, and its two $u$-lobes are disjoint and miss $1/3$ and $1$ as specified in
Proposition~\ref{prop:twists}(ii) (marked-fibre routing).
\item[(H11)] $a\le h_K<1/3$, $D=1-3h_K$, and
$h_K>\max_{G\in[-w_G,G_{\max}]}(O(G)-\mu\operatorname{slack}(G))$
(the common post-chart domain in \eqref{eq:p2-hK}).
\item[(H12)] $K_B$ has input radius $r_B=\sqrt{h_K/\pi}$ and satisfies the positive-parameter conditions and
\eqref{eq:K123} (the corner keyhole).
\item[(H13)] $\max(\pi d_B^2,\delta_{\mathrm{rel},B},\delta_{\mathrm{small},B})\le\delta_{\mathrm{kh}}$
(post-chart fibre containment).
\item[(H14)] $0<d_2<m$ and $\eta_1=\frac23-(2\delta_{\mathrm{kh}}+S_{\max})/D>0$, where
$S_{\max}=\max_{|c|\le m}\sqrt{c^2+q(c)^2}$ (positive simplex coordinates).
\item[(H15)] $m\le3\eta_1D$ (separation in the simplex transfer).
\item[Split] $2m<D-3\delta_{\mathrm{kh}}$ (disjoint post-chart bands before and after the corner keyhole).
\end{description}

Thus ``(H0)--(H15)'' groups the parameter conditions; the split inequality is included whenever all transfer
hypotheses are invoked. The program checks these conditions by margin guards, with tolerances for rounded versions
of exact identities such as (K1). The guards' count and floating-point tolerances are implementation details,
not additional mathematical hypotheses.
